% =========================================================================
% Kapitel 5 – Empirische Analyse nach CRISP-DM
% Entwurf (2026-07-18), Zahlen validiert gegen results/eignungspruefung/ und
% results/demo_modellierung/. Mit "% TODO" markierte Stellen nach dem finalen
% Trainingslauf (alle 3 Verfahren + Tuning) aktualisieren.
%
% -------------------------------------------------------------------------
% PREAMBLE-SETUP FÜR CODE-HIGHLIGHTING (einmalig in die Hauptdatei kopieren)
% -------------------------------------------------------------------------
% Empfehlung: Paket "listings" statt "minted". Begründung: minted benötigt
% -shell-escape und eine Python/Pygments-Installation auf dem Kompilier-System;
% das kollidiert häufig mit Uni-Templates, Overleaf-Kompatibilität und der von
% Schröter gewünschten PDF/A-Erzeugung. listings ist reines LaTeX, überall
% lauffähig und PDF/A-unkritisch.
%
% \usepackage{listings}
% \usepackage{xcolor}
% \definecolor{codegruen}{rgb}{0.0,0.45,0.0}
% \definecolor{codegrau}{rgb}{0.45,0.45,0.45}
% \definecolor{codeblau}{rgb}{0.10,0.20,0.65}
% \definecolor{codehintergrund}{rgb}{0.97,0.97,0.97}
% \lstdefinestyle{pythonstil}{
%   language=Python,
%   backgroundcolor=\color{codehintergrund},
%   basicstyle=\ttfamily\footnotesize,
%   keywordstyle=\color{codeblau}\bfseries,
%   commentstyle=\color{codegruen}\itshape,
%   stringstyle=\color{codegrau},
%   numbers=left, numberstyle=\tiny\color{codegrau}, numbersep=8pt,
%   frame=single, rulecolor=\color{codegrau},
%   breaklines=true, showstringspaces=false, captionpos=b,
%   literate={ä}{{\"a}}1 {ö}{{\"o}}1 {ü}{{\"u}}1 {ß}{{\ss}}1
%            {Ä}{{\"A}}1 {Ö}{{\"O}}1 {Ü}{{\"U}}1
% }
% \lstset{style=pythonstil}
% \renewcommand{\lstlistingname}{Quellcode}   % Beschriftung "Quellcode 5.1"
% Optional Verzeichnis: \lstlistoflistings nach dem Tabellenverzeichnis.
% -------------------------------------------------------------------------

\section{Empirische Analyse nach CRISP-DM}
\label{sec:empirie}

Die empirische Untersuchung folgt den Phasen des CRISP-DM-Prozessmodells.%
\footnote{Vgl. Wirth/Hipp (2000), S.~29\,ff.}
Entsprechend dem Schwerpunkt der Arbeit erhalten die Phasen Data Preparation,
Modeling und Evaluation das größte Gewicht; die Deployment-Phase ist nicht
Gegenstand der Untersuchung und wird in Abschnitt~6.3 als Limitation
eingeordnet.\footnote{Vgl. Schröer et al. (2021), S.~526\,ff.} Sämtliche
Analysen wurden in Python umgesetzt; die vollständigen Skripte liegen der
Arbeit bei, im Fließtext werden ausschließlich die methodisch entscheidenden
Ausschnitte als Beleg wiedergegeben.

% =========================================================================
\subsection{Business und Data Understanding}
\label{sec:business-data-understanding}

Ausgangspunkt der Analyse ist der in Kapitel~4 beschriebene integrierte
Datenbestand aus 720.258 Einsatzdatensätzen des San Francisco Fire Department
(2003--2026), sozioökonomischen Merkmalen des American Community Survey,
Kriminalitätskennzahlen des San Francisco Police Department sowie baulichen
Strukturmerkmalen der Parzellendatenbank Land Use 2020. Im Rahmen des Data
Understanding wurde eine systematische Eignungsprüfung durchgeführt, die
klärt, ob die drei zu vergleichenden Verfahren auf diesen Daten methodisch
zulässig und fair vergleichbar eingesetzt werden können.

Für die Zielgröße der Regression, die Anzahl der Einsätze je Stadtteil und
Monat, zeigt die explorative Analyse eine rechtsschiefe Zähldatenverteilung
(Mittelwert 63,4; Median 45; Maximum 451). Der Dispersionsindex von 61
(Varianz/Mittelwert) belegt eine ausgeprägte Überdispersion, sodass als
interpretierbare Count-Baseline eine Negative-Binomial-Regression anstelle
einer Poisson-Regression spezifiziert wird.%
\footnote{Vgl. Cameron/Trivedi (2013), S.~1\,ff.}

Für die Zielgröße der Klassifikation, die Einsatzart nach NFIRS-Systematik,
ergibt die Auszählung eine stark unbalancierte Verteilung: Fehlalarme stellen
mit 44,7\,\% die häufigste Kategorie, echte Brände lediglich 13,1\,\%. Die im
Exposé vorgesehene Vereinfachung auf eine binäre Klassifikation
(Brand vs.\ Nicht-Brand, 13,1/86,9\,\%) wird daher umgesetzt; dem
Klassenungleichgewicht wird durch Klassengewichte sowie die Gütemaße F1 und
AUROC anstelle der Trefferquote begegnet.

Zentraler Prüfpunkt ist die Zulässigkeit des linearen Verfahrens: Erst wenn
die grafische Analyse eine lineare Baseline erkennen lässt, darf ein lineares
Regressionsmodell eingesetzt werden. Die Streudiagramme der Prädiktoren gegen
die Zielgröße sowie eine OLS-Referenzschätzung zeigen, dass die
Stadtteilmerkmale bereits linear einen erheblichen Teil der Varianz erklären
(R\textsuperscript{2} $\approx$ 0{,}74; stärkster Einzelprädiktor ist der
Risiko-Gewerbe-Index mit $r = 0{,}69$). Eine lineare Baseline liegt somit vor,
Ridge Regression ist zulässig. Die Residuenanalyse offenbart jedoch auf der
Rohskala eine trichterförmige Heteroskedastizität sowie negative Vorhersagen
für Zählwerte; die Zielgröße wird für das lineare Modell daher
log-transformiert (Abschnitt~\ref{sec:feature-engineering}).

Die Prüfung auf Multikollinearität ergibt erhöhte, aber nicht extreme
Varianzinflationsfaktoren (Maximum: Medianeinkommen 8,8; Miete 7,7). Dieser
Befund begründet die Verfahrenswahl zusätzlich, da die L2-Regularisierung der
Ridge Regression genau für diese Datenlage entwickelt wurde;\footnote{Vgl.
Hoerl/Kennard (1970), S.~55\,ff.} für die baumbasierten Verfahren ist
Multikollinearität unkritisch, sie erschwert jedoch die Interpretation
einzelner Merkmalsbeiträge (Abschnitt~\ref{sec:shap}).

Schließlich zeigt die Autokorrelationsanalyse der Zielgröße eine für die
Untersuchung zentrale Eigenschaft: Mit einer Lag-1-Autokorrelation von 0,96
ist die monatliche Einsatzlast eines Stadtteils hochgradig persistent. Ein
naives Prognosemodell, das schlicht den Vormonatswert fortschreibt, erreicht
dadurch bereits R\textsuperscript{2} $=$ 0,88 und definiert damit die
maßgebliche Vergleichshürde für alle Machine-Learning-Verfahren
(Abschnitt~\ref{sec:evaluation}).

% =========================================================================
\subsection{Data Preparation}
\label{sec:data-preparation}

\subsubsection{Bereinigung und Behandlung fehlender Werte}
\label{sec:bereinigung}

Die Bereinigung setzt an drei Punkten an. Erstens enthalten die Quelldaten
269 mehrfach gemeldete Einsatznummern (0,04\,\%, davon 50 vollständig
identische Zeilen); diese wurden dedupliziert, da sie die Einsatzzählungen
systematisch überschätzen würden. Zweitens wurde die Ausrückzeit als Differenz
aus Ankunfts- und Alarmzeitpunkt berechnet und auf den plausiblen Bereich von
0 bis 60 Minuten begrenzt, wodurch rund 1,7\,\% der Datensätze entfielen
(negative Werte deuten auf Erfassungsfehler, extreme Werte auf
Protokollartefakte). Quellcode~\ref{lst:bereinigung} zeigt beide Schritte.

\begin{lstlisting}[caption={Deduplikation und Plausibilitätsfilter der
SFFD-Einsatzdaten (Ausschnitt aus \texttt{pipeline/02\_join.py})},
label=lst:bereinigung]
# Mehrfach gemeldete Einsatznummern aus den Quelldaten entfernen
df = df.drop_duplicates(subset=["incident_number"], keep="first")

# Ausrueckzeit berechnen und auf plausiblen Bereich begrenzen
df["response_time_min"] = (
    df["arrival_dttm"] - df["alarm_dttm"]
).dt.total_seconds() / 60
df = df[(df["response_time_min"] >= 0) & (df["response_time_min"] <= 60)]
\end{lstlisting}

Drittens wurde die Verknüpfung der Einsätze mit den jährlich erhobenen
ACS-Merkmalen strikt prognostisch gestaltet: Jeder Einsatz erhält den
\emph{letzten verfügbaren} ACS-Jahrgang (Erhebungsjahr $\leq$ Einsatzjahr).
Eine Zuordnung des zeitlich \emph{nächsten} Jahrgangs würde für frühe
Einsatzjahre Informationen aus der Zukunft einschleusen und damit eine
Datenleckage erzeugen, die die Prognosegüte systematisch überschätzt
(Quellcode~\ref{lst:acsjoin}).\footnote{Zur Leakage-Problematik bei
zeitstrukturierten Daten vgl. Bergmeir/Benítez (2012), S.~192\,f.}

\begin{lstlisting}[caption={Strikt prognostische Zuordnung des
ACS-Jahrgangs (Ausschnitt aus \texttt{pipeline/02\_join.py})},
label=lst:acsjoin]
# Jeder Einsatz erhaelt den letzten VERFUEGBAREN ACS-Snapshot
# (acs_jahr <= Einsatzjahr) - kein Blick in die Zukunft.
sffd["acs_year"] = sffd["year"].apply(
    lambda y: max([a for a in acs_years if a <= int(y)],
                  default=min(acs_years))
)
\end{lstlisting}

Nach demselben Prinzip werden die Kriminalitätsmerkmale je Stadtteil nur bis
zum Vorjahr des jeweiligen Einsatzes kumuliert. Die baulichen Merkmale liegen
ausschließlich als Querschnitt des Jahres 2020 vor und gehen daher als
zeitkonstante Strukturmerkmale ein; diese Einschränkung wird in
Abschnitt~6.3 diskutiert.

Bei den fehlenden Werten ist die Bildungsvariable maßgeblich: Die
ACS-Tabelle B15003 existiert im Jahrgang 2009 nicht, sodass die
Akademikerquote bei strikt prognostischer Zuordnung erst ab dem Einsatzjahr
2014 zur Verfügung steht. Anstelle einer Imputation über mehrere Jahre --
insbesondere anstelle eines Rückwärts-Auffüllens, das erneut
Zukunftsinformation einschleusen würde -- wird die Hauptanalyse auf den
Zeitraum ab 2014 begrenzt und der volle Zeitraum ohne dieses Merkmal als
Sensitivitätsanalyse berichtet. Der Stadtteil McLaren Park wird als
dokumentiertes Zensus-Artefakt ausgeschlossen (ausgewiesene Armutsquote von
0,90 bei nur 850 Einwohnern einer Parkfläche); Treasure Island und Mission
Bay erhalten erst ab dem ACS-Jahrgang 2021 Werte, Lakeshore weist keine
Baujahrangaben auf, sodass die betroffenen Stadtteil-Monate über den
NaN-Ausschluss entfallen.

\subsubsection{Räumlich-zeitliche Aggregation}
\label{sec:aggregation}

% ---- Übernahme aus Exposé (Fußnote KI-Verzeichnis Prompt 2) ----
Eine zentrale Entwurfsentscheidung betrifft die Granularität: Da die
Feature-Werte auf Stadtteilebene und jährlicher Erhebung vorliegen, würden
sich diese Werte bei einer Auswertung auf Einzeleinsatz-Ebene mehrfach
wiederholen und potenziell zu Pseudo-Signalen führen. Die Daten wurden daher
zeitlich auf Stadtteil~$\times$~Monat aggregiert; Zielgröße für die
Regression ist die Anzahl der Einsätze pro Stadtteil und
Monat.\footnote{In Anlehnung an Anthropic Claude (2026), KI-Verzeichnis
Prompt~2.}
% ---- Ende Übernahme ----

Gegenüber den denkbaren Alternativen ist diese Ebene wie folgt abzuwägen. Die
Einzeleinsatz-Ebene würde die Stichprobe künstlich auf über 700.000 Zeilen
aufblähen, obwohl die Prädiktoren nur auf Stadtteilebene variieren; die
Folge wären scheinbar hochsignifikante, tatsächlich aber redundante
Zusammenhänge. Eine Aggregation auf Census-Tract~$\times$~Jahr böte zwar
feinere räumliche Auflösung, scheitert jedoch daran, dass die Einsatzdaten
keine Tract-Zuordnung tragen, und würde mit der Jahresebene zugleich die
Saisonalität eliminieren. Die gewählte Ebene Stadtteil~$\times$~Monat ist
somit die kleinste Einheit, auf der Zielgröße und Prädiktoren konsistent
vorliegen und saisonale Dynamik erhalten bleibt.

Technisch wird ein vollständiges Raster aus Stadtteilen und Kalendermonaten
aufgespannt, sodass einsatzfreie Monate als echte Nullen in die Zählung
eingehen; der unvollständige Randmonat wird abgeschnitten. Fehlende
Merkmalswerte in rasterbedingt ergänzten Monaten werden ausschließlich
\emph{vorwärts} aus der Vergangenheit desselben Stadtteils aufgefüllt
(Quellcode~\ref{lst:aggregation}).

\begin{lstlisting}[caption={Aggregation auf Stadtteil~$\times$~Monat mit
vollständigem Raster (Ausschnitt aus
\texttt{modellierung/aggregation.py})}, label=lst:aggregation]
agg = (df.groupby(["stadtteil", "jahr", "monat"])
         .agg(anzahl_einsaetze=("einsatz_nummer", "count"),
              **{c: (c, "first") for c in PRAEDIKTOREN})
         .reset_index())

# Vollstaendiges Raster: Monate ohne Einsatz sind echte Nullen
raster = agg.set_index(["stadtteil", "jahr", "monat"]) \
            .reindex(idx).reset_index()
raster["anzahl_einsaetze"] = raster["anzahl_einsaetze"].fillna(0)

# Nur VORWAERTS fuellen - Rueckwaertsfuellen wuerde fehlende Werte
# mit Zukunftsinformation imputieren (Leakage).
raster[PRAEDIKTOREN] = (raster.groupby("stadtteil")[PRAEDIKTOREN]
                              .transform(lambda s: s.ffill()))
\end{lstlisting}

Die Zielgröße der Klassifikation verbleibt demgegenüber bewusst auf der
Einzeleinsatz-Ebene, da die Einsatzart ein Merkmal des einzelnen Ereignisses
ist und die einsatzbezogenen Zeitmerkmale (Stunde, Wochentag, Nacht) dort
ihre Erklärungskraft entfalten.

\subsubsection{Feature Engineering und finaler Analysedatensatz}
\label{sec:feature-engineering}

% ---- Übernahme aus Exposé ----
Im Rahmen des Feature-Engineerings wurden einzelne Prädiktoren aus den
Rohdaten neu berechnet: So wurden etwa der Anteil an Gebäuden mit Baujahr vor
1940 auf Stadtteilebene sowie ein Index für risikobehaftetes Gewerbe, auf
Basis der Anzahl von Industrie-, Gastronomie- und vergleichbaren Betrieben,
als eigenständige Merkmale konstruiert und dem Modell als zusätzliche
Prädiktoren bereitgestellt.
% ---- Ende Übernahme ----
Ergänzend wurden die sozioökonomischen Quoten (Armuts-, Akademiker- und
Leerstandsquote) als Zähler-Nenner-Verhältnisse im Intervall $[0,1]$
berechnet, die Mediane von Einkommen und Miete populationsgewichtet von der
Tract- auf die Stadtteilebene aggregiert und die Kriminalitätsanteile als
Gewalt- bzw.\ Eigentumsdelikte je Gesamtdelikte gebildet. Die Saisonalität
wird über eine zyklische Kodierung des Monats abgebildet, die den
künstlichen Sprung zwischen Dezember und Januar einer ordinalen Kodierung
vermeidet.

Über das Exposé hinaus wurden autoregressive Verlaufsmerkmale ergänzt. Den
Anlass gibt der Befund aus Abschnitt~\ref{sec:business-data-understanding}:
Bei einer Lag-1-Autokorrelation von 0,96 kann kein rein
struktur\-merkmalbasiertes Modell die naive Fortschreibung des
Vormonatswerts übertreffen. Je Stadtteil werden daher der Vormonatswert
(\texttt{lag\_1}), der Vorjahresmonat (\texttt{lag\_12}) und der gleitende
Dreimonatsmittelwert (\texttt{rolling\_mean\_3}) gebildet -- sämtlich streng
vergangenheitsbezogen (Quellcode~\ref{lst:lags}).

\begin{lstlisting}[caption={Streng vergangenheitsbezogene Lag-Merkmale je
Stadtteil (Ausschnitt aus \texttt{modellierung/demo\_modellierung.py})},
label=lst:lags]
g = daten.groupby("stadtteil")["anzahl_einsaetze"]
daten["lag_1"]          = g.shift(1)          # Vormonat
daten["lag_12"]         = g.shift(12)         # Vorjahresmonat
daten["rolling_mean_3"] = g.transform(        # Mittel der letzten
    lambda s: s.shift(1).rolling(3).mean())   # 3 Monate (ohne t)
\end{lstlisting}

Für die Auswertung werden zwei Merkmalsmengen unterschieden: Set~S umfasst
ausschließlich die Struktur- und Saisonmerkmale und dient der Beantwortung
der ersten Unterfrage nach dem Erklärungsbeitrag der Stadtteilmerkmale;
Set~S+L ergänzt die Lag-Merkmale und stellt die faire Prognoseaufgabe dar,
auf der die drei Verfahren verglichen werden. Modellspezifische
Transformationen erfolgen ausschließlich innerhalb der je Trainingsfenster
gefitteten Pipeline: Für die Ridge Regression werden alle Prädiktoren
z-standardisiert, die Zielgröße als $\log(1+y)$ modelliert und die
Lag-Merkmale ebenfalls $\log(1+x)$-transformiert, da rohe Zählwerte in einem
log-linearen Modell fehlspezifiziert sind; die baumbasierten Verfahren
trainieren unmittelbar auf den unveränderten Zählwerten. Für die
Klassifikation wird das Bataillon einheitlich für alle drei Verfahren
One-Hot-kodiert, sodass sämtliche Modelle eine identische Designmatrix
erhalten und Leistungsunterschiede ausschließlich algorithmisch bedingt sind.

Den finalen Analysedatensatz fasst Tabelle~\ref{tab:steckbrief} zusammen.
% TODO: Zahlen nach finalem Lauf ggf. aktualisieren (Stand: Demo 2026-07-18).

\begin{table}[htbp]
  \centering
  \caption{Steckbrief des finalen Analysedatensatzes}
  \label{tab:steckbrief}
  \begin{tabular}{ll}
    \hline
    \textbf{Merkmal} & \textbf{Ausprägung} \\
    \hline
    Regressionsebene & Stadtteil $\times$ Monat (vollständiges Raster) \\
    Hauptanalyse-Zeitraum & 2014--2026 (2015--2026 nach Lag-Anlauf) \\
    Beobachtungen (Regression) & 5.103 (39 Stadtteile $\times$ 133 Monate) \\
    Zielgröße Regression & Einsätze/Monat (Mittel 63,4; Max.\ 451) \\
    Prädiktoren & 11 Struktur + 2 Saison + 3 Lag \\
    Klassifikationsebene & Einzeleinsatz (719.989 Datensätze) \\
    Zielgröße Klassifikation & Brand vs.\ Nicht-Brand (13,1/86,9\,\%) \\
    Ausschlüsse & Duplikate; Ausrückzeit $\notin [0,60]$ min; \\
                & McLaren Park; Stadtteile ohne ACS-/Baujahrdaten \\
    \hline
  \end{tabular}
\end{table}

% =========================================================================
\subsection{Modellierung und Hyperparameter-Tuning}
\label{sec:modellierung}

\subsubsection{Validierungsdesign}

Da die Beobachtungen zeitlich geordnet sind, verbietet sich eine zufällige
Aufteilung in Trainings- und Testdaten: Sie würde Informationen aus der
Zukunft in das Training einschleusen und die Prognosegüte
überschätzen.\footnote{Vgl. Bergmeir/Benítez (2012), S.~192\,ff.} Die
Validierung erfolgt daher als Time-Series-Cross-Validation mit wachsendem
Trainingsfenster (expanding window): In drei Durchläufen wird jeweils auf
allen Monaten bis zum Zeitpunkt $t$ trainiert und auf den zwölf
Folgemonaten getestet (Quellcode~\ref{lst:tscv}). Alle Verfahren, Baselines
und beide Merkmalsmengen verwenden exakt dieselben Fold-Grenzen, sodass
Leistungsunterschiede ausschließlich auf die Verfahren zurückführbar sind.

\begin{lstlisting}[caption={Expanding-Window-Folds über Kalendermonate
(Ausschnitt aus \texttt{modellierung/demo\_modellierung.py})},
label=lst:tscv]
def zeit_folds(monate):
    """Training waechst, Testfenster (12 Monate) rollt nach vorn.
    Testmonate liegen stets NACH dem letzten Trainingsmonat."""
    folds = []
    for i in range(N_FOLDS):
        ende  = len(monate) - (N_FOLDS - 1 - i) * N_TEST_MONATE
        start = ende - N_TEST_MONATE
        folds.append((monate[:start], monate[start:ende]))
    return folds
\end{lstlisting}

\subsubsection{Verfahren und Referenzmodelle}

Verglichen werden Ridge Regression, Random Forest und XGBoost. Als
Referenzgrößen dienen erstens ein naives Modell, das je Stadtteil den
Vormonatswert fortschreibt, zweitens ein saisonaler Durchschnitt, der den
Mittelwert desselben Kalendermonats im Trainingsfenster verwendet, und
drittens die Negative-Binomial-Regression als interpretierbare
Count-Baseline. Für die Ridge Regression wird die in
Abschnitt~\ref{sec:feature-engineering} begründete log-lineare Spezifikation
als sklearn-Pipeline umgesetzt, deren Standardisierung ausschließlich auf dem
jeweiligen Trainingsfenster gefittet wird (Quellcode~\ref{lst:ridge}); die
Rücktransformation der Vorhersagen erfolgt vor der Bewertung, sodass alle
Gütemaße auf der Originalskala berechnet werden. In der Klassifikation
vertritt eine L2-regularisierte logistische Regression das lineare
Verfahren, da sie im Gegensatz zum Ridge-Klassifikator die für die
AUROC-Berechnung erforderlichen Wahrscheinlichkeiten liefert.

\begin{lstlisting}[caption={Log-lineare Ridge-Pipeline mit
trainingsfensterinterner Standardisierung}, label=lst:ridge]
modell = make_pipeline(StandardScaler(), Ridge(alpha=1.0))
modell.fit(X_train, np.log1p(y_train))     # Ziel: log(1+y)
y_hat = np.expm1(modell.predict(X_test))   # Bewertung: Originalskala
\end{lstlisting}

\subsubsection{Hyperparameter-Tuning}

Die Hyperparameter werden je Verfahren per Randomized Search
optimiert,\footnote{Vgl. Bergstra/Bengio (2012), S.~281\,ff.} wobei für
jedes Verfahren dasselbe Budget von 50 Iterationen gilt und die Bewertung
ausschließlich auf einem inneren Validierungsfenster erfolgt, das aus den
letzten zwölf Monaten des jeweiligen Trainingsfensters gebildet wird; die
Testmonate bleiben vom Tuning strikt unberührt. Die durchsuchten
Wertebereiche orientieren sich an der methodischen Literatur, für den Random
Forest insbesondere an den Empfehlungen von Probst et
al.\footnote{Vgl. Probst et al. (2019), S.~5\,ff.}
(Tabelle~\ref{tab:suchraeume}).

\begin{table}[htbp]
  \centering
  \caption{Suchräume des Randomized Search (50 Iterationen je Verfahren)}
  \label{tab:suchraeume}
  \begin{tabular}{ll}
    \hline
    \textbf{Verfahren} & \textbf{Suchraum} \\
    \hline
    Ridge & $\alpha \sim$ log-uniform $[10^{-3}, 10^{3}]$ \\
    Random Forest & \texttt{n\_estimators} 200--1000; \texttt{max\_depth}
      \{None, 10--40\}; \\
      & \texttt{min\_samples\_leaf} 1--20; \texttt{max\_features}
      \{sqrt, 0{,}3--1{,}0\} \\
    XGBoost & \texttt{n\_estimators} 200--2000; \texttt{learning\_rate}
      log-uniform $[0{,}005;\,0{,}3]$; \\
      & \texttt{max\_depth} 3--10; \texttt{subsample},
      \texttt{colsample\_bytree} 0{,}5--1{,}0; \\
      & \texttt{reg\_lambda} log-uniform $[10^{-2}, 10]$ \\
    Negative Binomial & kein Tuning (interpretierbare Referenz) \\
    \hline
  \end{tabular}
\end{table}

Sämtliche Läufe verwenden einen festen Startwert des Zufallsgenerators
(\texttt{random\_state} $=$ 42); Trainings- und Inferenzzeiten werden je
Verfahren und Fold mit \texttt{time.perf\_counter()} gemessen und als Median
aus drei Wiederholungen auf identischer Hardware berichtet.
% TODO: Hardware-Angabe ergänzen (CPU, RAM, Python-/Paketversionen).

% =========================================================================
\subsection{Evaluation und Modellvergleich}
\label{sec:evaluation}

Die Bewertung der Regressionsmodelle erfolgt anhand von RMSE, MAE und
R\textsuperscript{2}, jeweils je Fold sowie als Mittelwert über die Folds;
die Klassifikationsmodelle werden mit F1 (positive Klasse: Brand) und AUROC
bewertet. Tabelle~\ref{tab:ergebnisse-regression} zeigt den Stand der
Voruntersuchung mit untunierten Modellkonfigurationen.
% TODO: Nach dem finalen Lauf durch die getunten Ergebnisse aller drei
% Verfahren (inkl. XGBoost, NegBin-Baseline, Laufzeiten, Std.-Abw. je Fold)
% ersetzen und Klassifikationstabelle ergänzen.

\begin{table}[htbp]
  \centering
  \caption{Prognosegüte in der Time-Series-CV (3 Folds à 12 Testmonate;
  Voruntersuchung, untuniert; S $=$ Strukturmerkmale, L $=$ Lag-Merkmale)}
  \label{tab:ergebnisse-regression}
  \begin{tabular}{lrrr}
    \hline
    \textbf{Modell} & \textbf{RMSE} & \textbf{MAE} &
    \textbf{R\textsuperscript{2}} \\
    \hline
    Ridge (S+L)             & 21{,}6 & 12{,}0 & 0{,}91 \\
    Random Forest (S+L)     & 22{,}2 & 12{,}1 & 0{,}90 \\
    Naive Baseline (Vormonat) & 24{,}5 & 13{,}9 & 0{,}88 \\
    Random Forest (S)       & 24{,}6 & 13{,}5 & 0{,}88 \\
    Saisonaler Durchschnitt & 29{,}6 & 16{,}5 & 0{,}85 \\
    Ridge (S)               & 70{,}7 & 32{,}3 & 0{,}17 \\
    \hline
  \end{tabular}
\end{table}

Bereits die Voruntersuchung erlaubt drei Einordnungen, die für die
Interpretation des finalen Vergleichs leitend sind. Erstens definiert die
naive Fortschreibung des Vormonats -- nicht der saisonale Durchschnitt --
die relevante Vergleichshürde; ein Verfahren, das sie nicht übertrifft,
besitzt keinen praktischen Prognosemehrwert. Zweitens erklären die
Strukturmerkmale allein das Niveau der Stadtteile, nicht deren monatliche
Dynamik: Der Random Forest erreicht mit Set~S das Niveau der naiven
Baseline, indem er die Stadtteilniveaus über die Merkmalskombinationen
memoriert, während das lineare Modell mit Set~S deutlich zurückbleibt.
Drittens übertreffen beide Verfahrensklassen erst mit den
Verlaufsmerkmalen (Set~S+L) die naive Baseline. Der Modellvergleich auf
Set~S+L beantwortet damit die zweite Unterfrage, während der Kontrast
zwischen Set~S und Set~S+L gemeinsam mit der SHAP-Analyse
(Abschnitt~\ref{sec:shap}) den Erklärungsbeitrag der Strukturmerkmale
quantifiziert. Ergänzend werden die Ergebnisse je Fold ausgewiesen, da die
Prognoseaufgabe im jüngsten Testfenster erkennbar schwieriger ist als in den
älteren -- ein Hinweis auf Nichtstationarität, der in Abschnitt~6.3
aufgegriffen wird.

% =========================================================================
\subsection{Interpretierbarkeit und SHAP-Analyse}
\label{sec:shap}

Zur inhaltlichen Einordnung der Prognosemodelle werden SHAP-Werte
berechnet, die den Beitrag jedes Merkmals zu einer einzelnen Vorhersage als
additiv zerlegten Shapley-Wert ausweisen.\footnote{Vgl. Lundberg/Lee (2017),
S.~1\,ff.} Erklärt werden die jeweils final getunten Modelle des letzten
Folds; als Hintergrunddatensatz dient eine Zufallsstichprobe von 500
Beobachtungen ausschließlich aus dem Trainingsfenster, ausgewertet werden
die Testbeobachtungen (out-of-sample). Für die baumbasierten Verfahren kommt
der exakte \texttt{TreeExplainer} zum Einsatz, für das lineare Modell der
\texttt{LinearExplainer} auf den standardisierten Prädiktoren
(Quellcode~\ref{lst:shap}).

\begin{lstlisting}[caption={SHAP-Berechnung für das getunte
XGBoost-Modell (Ausschnitt aus \texttt{modellierung/shap\_analyse.py})},
label=lst:shap]
hintergrund = X_train.sample(500, random_state=42)  # nur Training!
explainer   = shap.TreeExplainer(
    modell, data=hintergrund,
    feature_perturbation="interventional")
shap_werte  = explainer(X_test)                     # out-of-sample
shap.plots.beeswarm(shap_werte)
\end{lstlisting}

Bei der Interpretation sind drei methodische Punkte zu beachten. Erstens
liegen die SHAP-Werte des log-linearen Ridge-Modells auf der
log-transformierten Skala, die der Baummodelle auf der Originalskala;
zwischen den Verfahren werden daher keine absoluten Beträge, sondern
Rangfolgen und normierte Anteile der mittleren absoluten SHAP-Werte
verglichen. Zweitens verteilen SHAP-Werte den Beitrag korrelierter Merkmale
auf die beteiligten Variablen; angesichts der festgestellten
Multikollinearität (Einkommen, Miete, Armutsquote) werden die Beiträge
zusätzlich zu Merkmalsgruppen aggregiert -- sozioökonomisch,
kriminalitätsbezogen, baulich, saisonal und autoregressiv --, deren
Anteilswerte die zentrale Auswertungsgröße zur ersten Unterfrage bilden.
% TODO: Gruppen-Balkendiagramm nach finalem Lauf einfügen (Abb. X).
Drittens erklären SHAP-Werte ausschließlich das Modellverhalten und erlauben
keine kausalen Aussagen über die Einsatzentstehung;\footnote{Vgl. Shmueli
(2010), S.~289\,ff.; Molnar (2022), Kap.~9.} die Formulierung der Befunde
beschränkt sich daher konsequent auf Prognosebeiträge und Assoziationen. Für
die Klassifikationsmodelle werden die SHAP-Werte auf der Margin-Skala
(Log-Odds) berechnet; die Koeffizienten der logistischen Regression dienen
als modellinterne Quervalidierung der SHAP-Rangfolge.

